% Model hierarchy diagram (fig:model_hierarchy).
% Each arrow points from a base model to an expanded model that nests it (the
% base is recovered by fixing the added parameters to zero); the edge label
% names the added parameters. Needs the TikZ libraries positioning and
% arrows.meta and the colours umlfill and umlborder (all in Setup/Settings.tex).
% Source diagram: docs/diagrams/08_model_hierarchy_beta.puml.
\begin{tikzpicture}[
    font=\fontsize{9}{11}\selectfont,
    model/.style={draw=umlborder, fill=umlfill, line width=0.6pt, rounded corners=4pt,
        minimum width=2.8cm, minimum height=0.8cm},
    nests/.style={umlborder, line width=0.7pt, -{Stealth[length=7pt]}},
    lbl/.style={text=black, inner sep=2pt},
    node distance=1.2cm and 1.0cm,
]
\node[model] (hmm) {HMM(1)};
\node[model, below left=of hmm] (ar) {AR(1)-HMM(1)};
\node[model, below right=of hmm] (beta) {HMM($\beta, 1$)};
\node[model, below right=of ar] (arbeta) {AR(1)-HMM($\beta, 1$)};

% Autoregressive emission first
\draw[nests] (hmm) -- node[lbl, above left] {$+\,\phi$} (ar);
\draw[nests] (ar) -- node[lbl, below left] {$+\,\beta$} (arbeta);

% Covariate-driven transition first
\draw[nests] (hmm) -- node[lbl, above right] {$+\,\beta$} (beta);
\draw[nests] (beta) -- node[lbl, below right] {$+\,\phi$} (arbeta);
\end{tikzpicture}
